While the academic prototype demonstrates significant technical achievements in perception, parallel execution, and Digital Twin synchronization, the transition from a laboratory environment to a commercial factory floor introduces rigorous economic and operational constraints. This chapter analyzes the industrial feasibility of deploying the proposed dual-arm collaborative cell, evaluating its cost-benefit ratio, integration hurdles with legacy systems, adherence to international safety standards, and impact on the human workforce.

\section{Economic Viability (Cost-Benefit Analysis)}

The economic justification for deploying an advanced, multi-robot assembly cell hinges on balancing the initial Capital Expenditure (CAPEX) against the projected Return on Investment (ROI) derived from increased throughput, reduced defect rates, and lowered Operational Expenditure (OPEX).

\subsection{Capital Expenditure (CAPEX)}
Deploying a Digital-Twin-enabled, dual-arm system requires a higher initial investment than traditional, single-arm automation. The primary CAPEX components include:
\begin{itemize}
    \item \textbf{Hardware:} The procurement of two industrial-grade manipulators. While this prototype utilizes a research-grade AR4 arm alongside the industrial ABB IRB120, a commercial deployment would likely scale to two industrial arms (e.g., dual ABB units or universal cobots) to guarantee continuous uptime and repeatability. Additional hardware costs include the high-resolution vision systems, IoT sensors, and the local edge-computing servers required to run the heavy ROS 2 and MoveIt 2 multithreading processes.
    \item \textbf{Software and Integration:} Development costs associated with mapping the physical cell to the RViz Digital Twin, configuring the MoveIt Task Constructor (MTC) pipelines, and setting up the cloud-bridge databases for remote control.
\end{itemize}

\subsection{Operational Expenditure (OPEX) and Targeted ROI}
Despite the higher CAPEX, the proposed architecture offers substantial OPEX reductions and unique ROI drivers compared to standard automation. As demonstrated in Chapter 7, the multithreaded parallel execution framework significantly minimizes arm starvation and idle time. Furthermore, the pre-execution validation within the Digital Twin reduces material scrap and machine downtime by preventing mechanical collisions before they occur physically.

Crucially, the economic viability of this specific architecture is highly dependent on the physical scale of the workpieces. While micro-electronics assembly poses significant challenges due to the limits of standard gripper precision, and heavy machinery exceeds the payload capacities of these arms, this system is optimized for \textbf{medium-scale manufacturing}. Target applications include:
\begin{itemize}
    \item Automotive sub-assembly kitting (e.g., dashboard components or fluid reservoirs).
    \item Consumer appliance internals (e.g., assembling motor housings or pump brackets).
    \item Medium-sized warehouse and logistics sorting, where asymmetric parts require dynamic grasp planning and collaborative handovers.
\end{itemize}
In these high-mix, medium-payload environments, the flexibility, high throughput, and zero-defect tolerances justify the advanced sensory and computational investments.